\documentclass[11pt]{article}
\usepackage{float}
\usepackage{amsmath}
\usepackage{booktabs}


\title{Where Does the Hedging Error Go? A Greek-by-Greek Decomposition for At-the-Money 0DTE Options under Heston}
\author{Harshith Ghanta}
\date{\today}

\begin{document}
\maketitle

\begin{abstract}
Zero Days to Expiration Options (0DTE) have increased in popularity since CBOE released daily SPY expirations in 2022. However, they are
difficult to hedge, and the sources of hedging error are not well understood. This paper takes the Heston model as a starting point
and uses an industry-standard Greek-by-Greek decomposition to measure how much each Greek moves the profit and loss of an
at-the-money delta hedge rebalanced at five-minute intervals. Shares are reported as gross absolute attribution shares computed
from model Greeks: because terms offset within a day, they measure the size of each term's contribution and not a signed split of
the realized hedging error. On that measure Gamma dominates, absorbing both diffusion convexity and any jump risk the model
cannot identify. On the days when the VIX travels furthest intraday, Vega's share rises from $1\%$ to $11\%$ while Gamma's falls
from $94\%$ to $82\%$ when $\epsilon$ = 0.5. The hedge also loses money persistently, by $14.4$ standard errors of the mean, because the
$\mathrm{VIX}^2$ variance fed to the model runs about three times the variance the underlying realizes.

\end{abstract}

\newpage
\section{Introduction}
This paper investigates the Greeks as sources of discrete hedging error for at-the-money 0DTE options under Heston.
It takes the standard compensated-Greek decomposition of a delta-hedged option's profit and loss and applies it to
at-the-money 0DTE calls, hedged at five-minute frequency on real SPY data. The contribution is a term-by-term measurement of how
much each Greek moves the hedged profit and loss, and how that ranking shifts with the volatility regime. The measure is a gross
absolute attribution share rather than a signed decomposition of the realized error, a distinction developed in
Section~\ref{sec:results}. Gamma dominates in every regime, but the balance between Gamma and Vega swings from roughly $94/1$ to
$82/11$ across terciles of intraday VIX movement, a sort whose endogeneity is discussed alongside the result.

\subsection{Black--Scholes Model and Greeks}
The Black--Scholes model, introduced by Fischer Black and Myron Scholes (1973)~\cite{blackscholes1973} and Robert Merton (1973)~\cite{merton1973},
provides a framework for pricing European options in a closed form that Heston expands on. The issue with Black--Scholes
is that it assumes constant volatility, which is not the case in real markets, especially for 0DTE options. The Heston
model, explained below, allows for stochastic volatility. Black--Scholes also introduced the concept of Greeks, which are
sensitivities of the option price to various parameters. They are listed below:
\begin{itemize}
    \item Delta ($\Delta$): Sensitivity of the option price to changes in the underlying asset price.
    \item Gamma ($\Gamma$): Sensitivity of Delta to changes in the underlying asset price.
    \item Vega ($\nu$): Sensitivity of the option price to changes in volatility.
    \item Theta ($\Theta$): Sensitivity of the option price to the passage of time.
    \item  Charm ($\frac{\partial \Delta}{\partial t}$): Sensitivity of Delta to the passage of time.
    \item Vanna ($\frac{\partial \Delta}{\partial \sigma}$): Sensitivity of Delta to changes in volatility.
    \item Volga, also called Vomma ($\frac{\partial \nu}{\partial \sigma}$): Sensitivity of Vega to changes in volatility.
    \item Veta ($\frac{\partial \nu}{\partial t}$): Sensitivity of Vega to the passage of time.
    \item Time convexity ($\frac{\partial^2 V}{\partial t^2}$): Second-order sensitivity of the option price to the passage of time.
\end{itemize}


\section{Background}
\subsection{Related Literature}


The decomposition in this paper relies heavily on various prior works. The Heston model \cite{heston1993}
is used as the underlying model for the option pricing and hedging found later in the paper. Along with that, the Euler--Maruyama
scheme \cite{Maruyama1955}, in the standard form given by Kloeden and Platen~\cite{kloedenplaten1992}, is used to simulate the
stochastic differential equations used in the implementation of the paper. The Greek decomposition is based on Bergomi's work \cite{bergomi2016}, which is used to decompose the hedging error into its
various components. Specifically, his ``break-even'' logic is used to attribute the hedging error to each Greek.
Bergomi's logic is the implied vs. realized volatility, where dominating volatility leads to price dominating the true price, and the Gamma variance drives the hedging error.
This implied-versus-realized logic was formalized by El~Karoui, Jeanblanc-Picqu\'e, and Shreve (1998)~\cite{elkaroui1998}, who show that the P\&L of a delta hedge run at the wrong volatility is the dollar-Gamma-weighted spread between the hedging and realized variances.

The result that discrete-hedging error scales as $\frac{1}{\sqrt{N}}$, where $N$ is the number of rebalances, originates with
Boyle and Emanuel (1980)~\cite{boyleemanuel1980}, who derive the variance of the hedging error of a discretely adjusted option
position. Kamal \cite{kamal1999} is the practitioner treatment of that result and the closest antecedent to this paper: a guide to
hedging errors when one cannot hedge continuously, written for a trading desk rather than for a journal. Bertsimas--Kogan--Lo
\cite{bertsimas2000} approach the same problem independently and from the continuous-time side,
finding that discrete tracking error grows with the granularity constant, $g$, which is built from accumulated dollar gamma. The formal contribution of Bertsimas et al. is
the expression $\frac{g}{\sqrt{N}}$; they also show that the kink in the call payoff does not break this convergence rate, but rather inflates $g$.
Dim, Eraker, and Vilkov (2023) \cite{dimeraker2023}
study how 0DTE order imbalances and the resulting dealer gamma exposure propagate into market volatility. That channel is microstructural; the attribution
here is at the contract level. No prior paper gives a term-by-term hedging-error attribution for 0DTE options under Heston.


\subsection{Choices and Rationale}
The Heston model was chosen for its stochastic volatility and easily implementable characteristic function. Also, the Heston parameters
are round, arbitrarily chosen values rather than calibrated estimates, except for $\eta$ (see the Methodology section for further information). Their specific combination violates the Feller condition,
so the variance process can reach zero; the simulation therefore uses a full-truncation Euler--Maruyama scheme~\cite{lord2010}, which floors the variance at
zero at each step.

\subsection{Heston Model Derivation}
The derivation here follows both Gatheral (2006)~\cite{gatheral2006} and Heston (1993)~\cite{heston1993}. The derivation starts with the Heston model's stochastic differential equations
for the underlying asset price $S$ and its variance $v$:
\[
dS = \mu S\,dt + \sqrt{v}\,S\,dX,
\qquad
dv = \kappa(\eta - v)\,dt + \epsilon \sqrt{v}\,dW,
\qquad
dX\,dW = \rho\,dt,
\]
There are two sources of randomness here, the price shock $dX$ and the variance shock $dW$, so the hedge takes two instruments: the stock, and a second option $V_1$ with its own exposure to the variance.
Take the portfolio, $\Pi$, with the option $V$ by shorting $\Delta S$, which is the
dollar value of the stock shorted. Short $\Delta_1$ shares of $V_1$ as well:
\[
\Pi = V - \Delta S - \Delta_1 V_1 .
\]
Applying It\^o's lemma to $V$ gives its change over an instant:
\[
dV = \left( \frac{\partial V}{\partial t}
+ \frac{1}{2} v S^2 \frac{\partial^2 V}{\partial S^2}
+ \rho \epsilon v S \frac{\partial^2 V}{\partial S\,\partial v}
+ \frac{1}{2} \epsilon^2 v \frac{\partial^2 V}{\partial v^2} \right) dt
+ \frac{\partial V}{\partial S}\,dS
+ \frac{\partial V}{\partial v}\,dv,
\]
Apply the same expansion to $V_1$, and substitute both into $d\Pi$ to obtain the following equation:
\begin{align*}
d\Pi &= \left( \frac{\partial V}{\partial t}
+ \frac{1}{2} v S^2 \frac{\partial^2 V}{\partial S^2}
+ \rho \epsilon v S \frac{\partial^2 V}{\partial S\,\partial v}
+ \frac{1}{2} \epsilon^2 v \frac{\partial^2 V}{\partial v^2} \right) dt \\
&\quad {} - \Delta_1 \left( \frac{\partial V_1}{\partial t}
+ \frac{1}{2} v S^2 \frac{\partial^2 V_1}{\partial S^2}
+ \rho \epsilon v S \frac{\partial^2 V_1}{\partial S\,\partial v}
+ \frac{1}{2} \epsilon^2 v \frac{\partial^2 V_1}{\partial v^2} \right) dt \\
&\quad {} + \left( \frac{\partial V}{\partial S} - \Delta_1 \frac{\partial V_1}{\partial S} - \Delta \right) dS
+ \left( \frac{\partial V}{\partial v} - \Delta_1 \frac{\partial V_1}{\partial v} \right) dv.
\end{align*}
The last line shows how to remove the risk: set $\Delta_1$ to cancel the $dv$ exposure, and $\Delta$ to cancel the $dS$ exposure,
\[
\Delta_1 = \frac{\partial V/\partial v}{\partial V_1/\partial v},
\qquad
\Delta = \frac{\partial V}{\partial S} - \Delta_1 \frac{\partial V_1}{\partial S},
\]
leaving a riskless portfolio. No-arbitrage then forces $\Pi$ to earn the risk-free rate, $d\Pi = r\Pi\,dt$. Because the same argument holds for any option, the terms involving $V$ and the terms involving $V_1$ must each equal a common function of $(S, v, t)$; following Heston, that function is written as $-\nu_{\mathrm{Q}}$:
\[
\begin{aligned}
\frac{1}{\partial V/\partial v} \biggl(
&\frac{\partial V}{\partial t}
+ \frac{1}{2} v S^2 \frac{\partial^2 V}{\partial S^2}
+ \rho \epsilon v S \frac{\partial^2 V}{\partial S\,\partial v} \\
&{} + \frac{1}{2} \epsilon^2 v \frac{\partial^2 V}{\partial v^2}
+ rS \frac{\partial V}{\partial S}
- rV \biggr) = -\nu_{\mathrm{Q}},
\end{aligned}
\]
where
\[
\nu_{\mathrm{Q}} \equiv \kappa(\eta - v) - \lambda v
\]
is the risk-neutral drift of the variance, with $\lambda$ the market price of volatility risk. Rearranging gives the Heston pricing PDE:
\[
\begin{aligned}
&\frac{\partial V}{\partial t}
+ \frac{1}{2} v S^2 \frac{\partial^2 V}{\partial S^2}
+ \rho \epsilon v S \frac{\partial^2 V}{\partial S\,\partial v}
+ \frac{1}{2} \epsilon^2 v \frac{\partial^2 V}{\partial v^2} \\
&\quad {} + rS \frac{\partial V}{\partial S}
+ \big[\kappa(\eta - v) - \lambda v\big] \frac{\partial V}{\partial v}
- rV = 0.
\end{aligned}
\]

\section{The Error Model}
\subsection{Derivation}

The Error model is derived from the Heston model; the two coincide in the continuous limit and diverge as the hedge becomes discrete.
The Heston model assumes continuous hedging; therefore, $\Delta t \to 0$. However, in practice, hedging is done at discrete intervals,
 which introduces hedging error. To decompose said error, the following model is proposed based on Bergomi's compensated-Greek decomposition, adding the discrete-time terms (charm, veta, and time convexity) that
 arise from expanding the profit and loss to second order in $\Delta t$~\cite{bergomi2016}:
\[
\Delta\Pi = \Delta V - \frac{\partial V}{\partial S}\,\Delta S
- r\,\Big( V - \frac{\partial V}{\partial S}\,S \Big)\,\Delta t .
\]
\begin{align*}
\Delta\Pi &= \frac{\partial V}{\partial t}\,\Delta t
+ \frac{\partial V}{\partial S}\,\Delta S
+ \frac{\partial V}{\partial v}\,\Delta v \\
&\quad {} + \frac{1}{2} \frac{\partial^2 V}{\partial S^2}\,(\Delta S)^2
+ \frac{1}{2} \frac{\partial^2 V}{\partial v^2}\,(\Delta v)^2
+ \frac{\partial^2 V}{\partial S\,\partial v}\,\Delta S\,\Delta v \\
&\quad {} + \frac{\partial^2 V}{\partial S\,\partial t}\,\Delta S\,\Delta t
+ \frac{\partial^2 V}{\partial v\,\partial t}\,\Delta v\,\Delta t
+ \frac{1}{2} \frac{\partial^2 V}{\partial t^2}\,(\Delta t)^2 \\
&\quad {} - \frac{\partial V}{\partial S}\,\Delta S
- r\,\Big( V - \frac{\partial V}{\partial S}\,S \Big)\,\Delta t .
\end{align*}
\[
\frac{\partial V}{\partial t} = rV - rS\,\frac{\partial V}{\partial S}
- \frac{1}{2} v S^2\,\frac{\partial^2 V}{\partial S^2}
- \rho \epsilon v S\,\frac{\partial^2 V}{\partial S\,\partial v}
- \frac{1}{2} \epsilon^2 v\,\frac{\partial^2 V}{\partial v^2}
- \nu_{\mathrm{Q}}\,\frac{\partial V}{\partial v} ,
\]
\begin{align*}
\Delta\Pi &= \frac{1}{2} \frac{\partial^2 V}{\partial S^2}\big[ (\Delta S)^2 - v S^2\,\Delta t \big]
+ \frac{1}{2} \frac{\partial^2 V}{\partial v^2}\big[ (\Delta v)^2 - \epsilon^2 v\,\Delta t \big] \\
&\quad {} + \frac{\partial^2 V}{\partial S\,\partial v}\big[ \Delta S\,\Delta v - \rho \epsilon v S\,\Delta t \big]
+ \frac{\partial V}{\partial v}\big[ \Delta v - \nu_{\mathrm{Q}}\,\Delta t \big] \\
&\quad {} + \frac{\partial^2 V}{\partial S\,\partial t}\,\Delta S\,\Delta t
+ \frac{\partial^2 V}{\partial v\,\partial t}\,\Delta v\,\Delta t
+ \frac{1}{2} \frac{\partial^2 V}{\partial t^2}\,(\Delta t)^2 .
\end{align*}
Not every term above is discretization error. Summed over a day, the Gamma, Vanna, and Volga
compensators vanish as the hedging interval shrinks, since realized quadratic variation converges to the model's $\int v S^2\,dt$ and
likewise for the cross and variance brackets, while the charm, veta, and time-convexity terms are of order $\Delta t$ or smaller.
Those six terms measure finite-mesh error: the price of rebalancing every five minutes instead of continuously. The Vega term
behaves otherwise. Its sum converges to $\int \frac{\partial V}{\partial v}\,\epsilon\sqrt{v}\,dW$, a random quantity that
survives any rebalancing frequency. That is the variance risk a delta-only hedge leaves open, since the hedge holds no second
option. The Heston derivation above removed this risk with the hedging option $V_1$; the experiment below deliberately holds only
the stock, as most 0DTE hedgers do, so the Vega term tracks unhedged risk rather than finite-mesh error.

\subsection{Methodology}
\label{sec:methodology}
The Heston parameters were not calibrated to data: apart from $\eta$, they are round, arbitrarily chosen values.
$\eta$ was set from the VIX-implied variance in the sample. Given that information, the results in the
next section are not to be taken as absolute, but rather as a proof of concept. The parameters used are as follows:
\begin{itemize}
    \item $\kappa = 3.0$ (rate of mean reversion)
    \item $\eta = 0.04$ (long-run variance; the sample average of $(\mathrm{VIX}/100)^2$ is $0.035$, rounded)
    \item $\epsilon = 0.5$ (volatility of volatility)
    \item $\rho = -0.7$ (correlation between asset and variance)
    \item $\lambda = 0$ (market price of volatility risk)
    \item $r = 0.045$ (risk-free rate, used in the financing term of the hedged P\&L)
\end{itemize}

Two of these are uncalibrated free choices worth probing, $\kappa$ and $\epsilon$, and
Section~\ref{sec:results} reports a sweep over both. The mean-reversion rate $\kappa$ is expected to be
close to inert on this horizon: it enters only through $\nu_{\mathrm{Q}}\,\Delta t = \kappa(\eta - v)\,\Delta t$, which for
$\kappa = 3$ and a five-minute step is of order $10^{-6}$, against a typical per-bar $\Delta v$ two orders
of magnitude larger. The vol-of-vol $\epsilon$ has no such protection, since it scales the Volga
compensator as $\epsilon^2 v\,\Delta t$ and the Vanna compensator as $\rho\epsilon v S\,\Delta t$.

Five-minute SPY data was used as a proxy for the underlying asset; the pipeline rescales each day to open at 1, so the results are invariant to the units of the price series. The sample covers 56 trading days, from March 31st, 2026 through June 18th, 2026.
That window contains every regular NYSE session in the range apart from Good Friday and Memorial Day, which is where the count of 56 comes from.
The option prices were computed using the Heston model, and the hedging error was measured by rebalancing the delta hedge at each five-minute bar.
The contributions of each Greek to the overall hedged profit and loss were then analyzed and their results are found
in the following section. Better data would use actual option pricing instead of proxies, but option data is expensive and currently a limitation
for this paper.

In terms of the derivation itself, the specific terms above were chosen to allow
for all the terms to be expressed as Greeks, and have a discrete and nondiscrete component. The discrete component is the difference between the actual change
in the underlying asset and the expected change, while the nondiscrete component is the expected change itself. The discrete component, following Bergomi's 2016 book,
is what contributes to the
hedging error, while the nondiscrete component is what is expected to be hedged away. If the limit is taken, the model collapses back to the Heston model for every term except Vega:
the other compensators go to zero with the hedging interval, while the Vega term converges to the unhedged variance risk described above.

In the implementation, the trapezoid rule was used to compute the integrals in the Heston pricer, and the Greeks were computed using finite differences.
The code is available in the accompanying Python file on GitHub. The $\omega$ grid begins at $10^{-8}$ to avoid the $1/(i\omega)$ singularity at zero, and the upper limit of
the integral is set to 2500 to ensure convergence for 0DTE options. The number of points in the grid is set to 2048, which is enough for the reported prices to
stabilize. If using a stronger computer, the number of points can be increased to improve accuracy, but at the cost of longer computation time. With testing, 
it has been found that if $\lambda$ is set to 0, then the results are not sensitive to the number of points in the grid, as the integrals converge quickly. 
However, if $\lambda$ is nonzero, then the results can be sensitive to the number of points in the grid, as the integrals converge more slowly. Also, setting $\lambda$ to a nonzero value
can introduce numerical instability in the integrals, as the integrands can become highly oscillatory, therefore making the pricer and the vega compensator disagree.

Two distinct experiments appear below, and the distinction governs how each residual should be read. The empirical study of
Section~\ref{sec:results} is not simulated at all: $\Delta S$ is the realized five-minute change in the observed underlying and
$\Delta v$ is the realized change in the squared VIX. The validation of Section~\ref{sec:validation} is simulated, using the
Euler--Maruyama scheme to generate paths from the Heston dynamics.

On the simulated paths the price is carried in levels rather than log prices. An Euler step
$\Delta S = rS\,\Delta t + \sqrt{v}\,S\,\Delta X$ has conditional variance $vS^2\,\Delta t$, which is exactly the quantity the
Gamma compensator $(\Delta S)^2 - vS^2\,\Delta t$ subtracts. Because the paths are advanced in several sub-steps per hedge bar,
simulating in levels keeps that compensator mean-zero up to an $O(\Delta t^2)$ correction, of the same order as the Taylor
truncation itself, so any residual left in the \emph{simulated} test reflects second-order error rather than a mismatch
between how the paths were generated and how the error was decomposed. A log-Euler step would reintroduce such a mismatch through
the $-\frac{1}{2}v$ Itô correction, which the level scheme does not carry.

No such guarantee is available on real data, and the gap is large. There $\Delta S$ is whatever the market did, while
$vS^2\,\Delta t$ is built from a VIX-derived variance that was never fitted to the realized path. Summed within each day, the
realized quadratic variation $\sum (\Delta S)^2$ averages $4.38 \times 10^{-5}$ across the sample, against
$1.41 \times 10^{-4}$ for the model's $\sum v S^2\,\Delta t$. The model quantity is roughly three times the realized one, and it
exceeds it on 53 of the 56 days. The Gamma compensator is therefore not mean-zero on the empirical sample but reliably negative,
and because an at-the-money call is long Gamma the decomposition's largest term carries a systematic negative drift. The residual
reported in Section~\ref{sec:results} accordingly mixes second-order Taylor truncation with error in the variance proxy, and the
two residuals in this paper measure different things and should not be compared directly.


The variance path is read off the VIX, which is a proxy in two respects. First, $\mathrm{VIX}^2$ is a
risk-neutral quantity rather than a physical one, and it embeds a variance risk premium that the model,
with $\lambda = 0$, does not capture. Second, and more consequentially,
$\mathrm{VIX}^2$ is the expected \emph{average} variance over the coming thirty days, whereas the model
calls for the instantaneous spot variance $v_t$. Because variance mean-reverts, a thirty-day forward
average is pulled toward the long-run level and barely moves intraday, so the variance terms are only
lightly exercised on real data. Cboe's
VIX1D index, introduced in 2023 and computed from zero- to one-day SPY options~\cite{cboe2023vix1d}, is the horizon-matched
measure and should replace the VIX in future work.

This paper does not perform the VIX1D substitution. The expected effect is directional and stateable: a
horizon-matched variance proxy would move intraday, activating the Vega, Vanna, and Volga terms that the
thirty-day VIX suppresses, and would remove the roughly threefold gap between model and realized variance,
shrinking the systematic negative drift in the Gamma compensator. Both effects push in the same direction:
Gamma's share should fall from the level reported here, and the drift documented in
Section~\ref{sec:results} should attenuate or vanish. The rerun is a test of this paper's central number,
and it is the first thing that should be done.

\subsection{Validation}
\label{sec:validation}
According to Bertsimas, Kogan, and Lo (2000)~\cite{bertsimas2000}, the granularity constant $g$, which measures how difficult an option is to hedge
in discrete time and sets the size of the tracking error, blows up in the terminal interval. The Taylor expansion also
fails as the option approaches expiry, which is one candidate explanation for the 9\% residual reported in the next subsection.

The decomposition algebra itself was checked away from that corner. A 60-day option was simulated under the Heston engine and
hedged daily, with the position closed 20 days before expiry so that no step straddles the payoff kink and every Greek stays
smooth; there the seven terms reproduced the simulated hedged P\&L to within 2.35\%, at a correlation of 0.9997. That is a test of
whether the expansion is correct, and it passes.

The terminal-interval explanation was then tested directly, and on the real data rather than by analogy to the 60-day run. The
same 56 days were hedged the same way, but the position was closed a number of bars before the bell and marked out at the pricer
value instead of being held into expiry; every other element of the experiment is unchanged, so distance from expiry is the only
thing that varies. The residual falls monotonically and steeply as the position is closed earlier: $9.0\%$ held to expiry,
$8.1\%$ closed one five-minute bar early, $7.7\%$ two bars early, $7.2\%$ three bars early, $5.8\%$ half an hour early, $4.7\%$
one hour early, and $4.0\%$ two hours early. Closing merely one bar before the bell removes almost a percentage point of residual,
and closing thirty minutes early removes more than a third of it. The residual is therefore concentrated in the final intervals,
exactly as the granularity constant blowing up in the terminal interval would predict. This is a controlled version of the
comparison the 60-day run could only gesture at: the maturity, hedge frequency, and data source are all held fixed, so the fall
can be attributed to distance from expiry alone.

The 60-day simulation of the previous paragraph remains useful as an algebra check but is not the terminal-interval test: it
differs from the 0DTE study in maturity, hedge frequency, and data source at once, and by the argument of
Section~\ref{sec:methodology} its residual is not even the same quantity, being second-order truncation on paths where the Gamma
compensator is mean-zero to leading order, whereas the empirical residual also contains variance-proxy error. The direct truncation
test above is the controlled statement; the 60-day run is corroborating context.



\subsection{Results and Testing}
\label{sec:results}
The decomposition was tested on 56 full trading days of five-minute data (78 bars per day, so 77 hedge intervals). Each day is rescaled to open at 1, the variance path is read off the squared VIX, and the option is a fresh at-the-money 0DTE call re-formed each morning and delta-hedged at every bar. Summing the seven compensated terms reproduces the realized hedged P\&L to within a $9\%$ residual, defined as the mean absolute unexplained P\&L across the 56 days divided by the mean absolute daily P\&L. The mean daily hedged P\&L is $-0.00235$ in the units where the underlying opens at 1, with a standard deviation across days of $0.00122$.

That mean is not a sampling artifact. With a standard deviation of $0.00122$ across $D = 56$ days, the standard error of the mean is $0.00016$, placing the average daily loss $14.4$ standard errors below zero; 53 of the 56 days lose money. A delta hedge run on this data loses persistently rather than occasionally. The source is the variance input, measured in Section~\ref{sec:methodology}: $\mathrm{VIX}^2$ is a risk-neutral quantity carrying a variance risk premium, while $\lambda = 0$ gives the model a physical variance drift, so the $vS^2\,\Delta t$ subtracted inside the Gamma compensator runs about three times the variance the underlying actually realized and exceeds it on the same 53 days. An at-the-money call is long Gamma, so the position pays that spread on almost every day in the sample. The drift is the variance risk premium, arriving through the Gamma term rather than through Vega.

Table~\ref{tab:attribution} reports each term's size. The percentages are gross absolute model-attribution shares: each term's average absolute daily contribution as a fraction of the sum across all seven terms, computed from model Greeks. This measure deserves emphasis, because it is weaker than the phrase ``where the error goes'' suggests. Terms can and do offset within a day, so a share here says how much a term \emph{moves}, and not how much of the realized hedging error it accounts for. A signed decomposition, in which the seven contributions are attributed to the realized error with their signs intact and made to sum to it, would support the stronger claim; this paper does not construct one. In the table, Vega, Vanna, Volga, and Veta denote derivatives with respect to the variance $v$, not the volatility $\sigma$ of Section~1; the two are related by chain rule but are not numerically identical. The overall shares carry a percentile-bootstrap 95\% interval over days ($B = 10{,}000$ resamples of the 56 days, shares recomputed on each). The shares are reported first over all days, then split into three regimes by how far the VIX travelled intraday, measured in VIX index points.

\begin{table}[H]
\centering
\setlength{\tabcolsep}{4pt}
\begin{tabular}{lrcrrr}
\toprule
Term & Overall & 95\% CI & Calm & Mid & Turbulent \\
\midrule
Gamma $\left(\frac{1}{2}\,\partial^2 V/\partial S^2\right)$          & 89.7 & [88.0,\,91.2] & 93.6 & 91.5 & 82.1 \\
Vega $\left(\partial V/\partial v\right)$                            &  4.8 & [\phantom{0}3.6,\,\phantom{0}6.2] &  1.2 &  3.5 & 11.3 \\
Vanna $\left(\partial^2 V/\partial S\,\partial v\right)$             &  2.3 & [\phantom{0}1.9,\,\phantom{0}2.9] &  2.2 &  2.0 &  3.0 \\
Charm $\left(\partial^2 V/\partial S\,\partial t\right)$             &  1.5 & [\phantom{0}1.1,\,\phantom{0}1.9] &  1.5 &  1.3 &  1.7 \\
Time convexity $\left(\frac{1}{2}\,\partial^2 V/\partial t^2\right)$ &  1.3 & [\phantom{0}0.9,\,\phantom{0}1.7] &  1.2 &  1.3 &  1.4 \\
Volga $\left(\frac{1}{2}\,\partial^2 V/\partial v^2\right)$          &  0.2 & [\phantom{0}0.2,\,\phantom{0}0.3] &  0.3 &  0.2 &  0.2 \\
Veta $\left(\partial^2 V/\partial v\,\partial t\right)$              &  0.1 & [\phantom{0}0.1,\,\phantom{0}0.1] &  0.0 &  0.1 &  0.2 \\
\bottomrule
\end{tabular}
\caption{Gross absolute model-attribution share of each decomposition term, in percent, as a fraction of the summed average absolute daily contributions. The interval is a percentile bootstrap over days ($B = 10{,}000$ resamples of the 56 days). Days are sorted into terciles by the absolute intraday change in the VIX, measured in VIX index points: calm ($n=19$, mean absolute change $0.12$ points), mid ($n=18$, $0.47$ points), and turbulent ($n=19$, $1.57$ points). Shares are rounded to one decimal, so a Veta share of $0.0$ in the calm bucket means below $0.05$, not zero.}
\label{tab:attribution}
\end{table}

Gamma dominates in every regime, as one would expect of an option hedged only in the underlying: on a typical day roughly nine-tenths of the total absolute term movement sits in the $\frac{1}{2}\,\partial^2 V/\partial S^2$ term. The regime changes where the remaining tenth sits. On calm days Gamma carries $94\%$ and Vega barely $1\%$; on turbulent days Gamma's share falls to $82\%$ while Vega climbs to $11\%$, taking up almost all of the share Gamma gives up. Vanna, Charm, time convexity, Volga, and Veta stay within a few percent throughout, and by the bootstrap intervals in Table~\ref{tab:attribution} the Charm and time-convexity shares are too close to rank against each other.

The regime pattern needs one caveat. The Vega term is $\partial V/\partial v\,[\Delta v - \nu_{\mathrm{Q}}\Delta t]$, so it is larger by construction on days when the variance proxy moved; sorting days by exactly that quantity and then observing that Vega's share rises is partly circular. To check that the swing is a feature of turbulent markets rather than of the sort, the same 56 days were re-sorted into terciles by the absolute intraday move in the underlying, $|\Delta S|$, a criterion that does not appear in the Vega term. The swing survives: Gamma runs $91.8\%$ to $84.1\%$ from the calmest to the most active tercile, and Vega $2.4\%$ to $9.9\%$. The effect is attenuated relative to the $|\Delta\mathrm{VIX}|$ sort but clearly present, so it is not an artifact of the sorting variable. The two sorts are not independent, correlating at $0.80$ on this sample, so this is a partial control rather than a clean one.

The two remaining robustness questions are the mean-reversion rate $\kappa$ and the vol-of-vol $\epsilon$, both chosen rather than calibrated. Each was swept over three values with all other parameters and the data held fixed; Table~\ref{tab:sensitivity} reports both. The two behave very differently. Varying $\kappa$ across $\{1, 3, 5\}$ moves no share by as much as a tenth of a percentage point: the overall Gamma share is $89.7\%$ in all three cases and the turbulent-tercile split is unchanged. This confirms numerically what the magnitude argument of Section~\ref{sec:methodology} anticipated, that on a five-minute 0DTE horizon $\kappa$ enters only through the negligible $\nu_{\mathrm{Q}}\,\Delta t$ term and cannot affect the attribution. Varying $\epsilon$ across $\{0.5, 1.0, 2.0\}$ is a different story. Gamma remains the largest term by a wide margin, but its overall share is genuinely $\epsilon$-dependent, falling from $89.7\%$ at $\epsilon = 0.5$ to $79.2\%$ at $\epsilon = 2.0$. The share it sheds goes to Vanna and Volga, the cross- and second-variance terms that scale in $\epsilon$, and not to Vega, which stays near $4.5\%$ regardless. The regime pattern of Table~\ref{tab:attribution} is intact throughout: Gamma falls from the calm to the turbulent tercile and Vega rises to roughly 10 to 11 percent in the turbulent bucket in every row. The reconstruction residual is essentially unchanged by either parameter, at about $9\%$. Because calibrated short-horizon index vol-of-vol is commonly well above $0.5$, the honest reading of the headline is that Gamma carries roughly four-fifths to nine-tenths of the term movement, with the precise figure set by $\epsilon$ and not by $\kappa$.

\begin{table}[H]
\centering
\begin{tabular}{lrrrrrrr}
\toprule
& \multicolumn{4}{c}{Overall share (\%)} & \multicolumn{2}{c}{Turbulent (\%)} & Resid. \\
\cmidrule(lr){2-5}\cmidrule(lr){6-7}
& Gamma & Vega & Vanna & Volga & Gamma & Vega & (\%) \\
\midrule
\multicolumn{8}{l}{\emph{Vary mean reversion $\kappa$ (with $\epsilon = 0.5$):}} \\
\quad $\kappa = 1$   & 89.7 & 4.8 & 2.4 & 0.2 & 82.1 & 11.4 & 9.0 \\
\quad $\kappa = 3$   & 89.7 & 4.8 & 2.3 & 0.2 & 82.1 & 11.3 & 9.0 \\
\quad $\kappa = 5$   & 89.7 & 4.8 & 2.3 & 0.2 & 82.2 & 11.3 & 9.0 \\
\midrule
\multicolumn{8}{l}{\emph{Vary vol-of-vol $\epsilon$ (with $\kappa = 3$):}} \\
\quad $\epsilon = 0.5$ & 89.7 & 4.8 & 2.3 & 0.2 & 82.1 & 11.3 & 9.0 \\
\quad $\epsilon = 1.0$ & 86.7 & 4.6 & 4.7 & 1.2 & 78.9 & 10.9 & 8.9 \\
\quad $\epsilon = 2.0$ & 79.2 & 4.3 & 9.3 & 4.7 & 71.4 & 10.0 & 8.9 \\
\bottomrule
\end{tabular}
\caption{Attribution shares as $\kappa$ and $\epsilon$ are each swept, with the other parameters and the data held fixed. $\kappa$ is inert: all three rows are identical to within a tenth of a percentage point. $\epsilon$ matters: Gamma cedes share to Vanna and Volga as it grows, while Vega is unmoved. The regime split (Gamma falling and Vega near 10 to 11 percent in the turbulent tercile) holds in every row, and the reconstruction residual stays near $9\%$ throughout. The $\kappa = 3$, $\epsilon = 0.5$ baseline is shared by both panels.}
\label{tab:sensitivity}
\end{table}

\subsection{Limitations and Future Work}
Four limits on the scope of these results should be stated plainly. The option studied is at the money at inception on every one
of the 56 days, so the attribution is an at-the-money result and does not carry to the out-of-the-money strikes that make up much
of 0DTE volume. The at-the-money case was chosen deliberately: it is where dealer gamma exposure concentrates~\cite{dimeraker2023},
where the hedging problem is most acute, and where the Gamma-dominance result is most cleanly interpretable. Out-of-the-money
strikes, which carry proportionally more Vega and Vanna, are left to future work. The parameters are a single uncalibrated set;
the $\kappa$ and $\epsilon$ sweeps of Section~\ref{sec:results} show Gamma stays dominant throughout but that its precise share depends on $\epsilon$,
and no sweep over $\eta$, $\rho$, or $\lambda$ was attempted. The variance path is a thirty-day risk-neutral proxy standing in for an
instantaneous physical variance, and it runs about three times the realized intraday variance, which flatters the Gamma share
rather than merely muting the variance terms. Finally, the attribution metric is a gross absolute share, so it measures the size
of each term's movement and not a signed split of the realized error.

The model is limited by the accuracy of the Heston parameters, an issue one can easily fix with better data. Another limitation is found in the premise itself:
all Greeks are unhedgeable in the terminal interval, because there exists no rebalance prior to expiry. Given that, market makers and retail traders alike should be
wary of 0DTE options, which are risky and difficult to hedge, especially as they near expiration.
This paper locates the hedging error and does not provide a definitive solution to the problem. To make this model better, a Milstein scheme \cite{gatheral2006} could be used to simulate the variance process, which would provide a more accurate
representation of the underlying asset's volatility.

A further limitation is structural. Heston is a pure-diffusion model, so the decomposition carries no jump term; any intraday jump in the underlying enters only
through $(\Delta S)^2$ and is therefore absorbed into the Gamma term, indistinguishable from ordinary diffusive movement. For 0DTE the omission is costly.
Bo\v{z}ovi\'c (2025)~\cite{bozovic2025}, working with five-minute SPY options, estimates the implied jump-risk premium at nearly twice the combined diffusion and volatility premia,
with jump intensity clustering at the open and close. The roughly ninety-percent Gamma share reported above should therefore be read as diffusion and jump risk combined;
isolating the two would require adding a jump term to the decomposition.


\section{Conclusion}
Learning the different attributions of the error is important for understanding the risks of 0DTE options and
using them effectively. On a gross absolute attribution measure, the majority of the movement is carried by Gamma, which the
model conflates with jumps. The remainder is mostly regime-dependent: Vega, the variance risk a delta-only hedge cannot remove at any rebalancing frequency, takes a larger share as the regime turns turbulent, and that swing survives a re-sort on the underlying's own move rather than the VIX.
The other Greeks, while not negligible, are relatively small and do not change much across regimes.
These results are preliminary. They describe at-the-money calls at a single baseline parameter set, robust to the mean-reversion rate but not to the vol-of-vol, with a thirty-day variance proxy that runs about three times realized variance, and further work is needed before the attribution can be claimed for 0DTE hedging in general.
Future work should focus on improving the Heston parameter estimates, adding a jump term to the decomposition, and testing the model on a wider range of data.


\appendix
\section{Glossary of Variables}

The sensitivity Greeks ($\Delta$, $\Gamma$, $\nu$, $\Theta$, and the higher-order terms) are
defined in Section~1. The symbols below collect the model, market, and hedging notation that the
Heston and derivation sections introduce along the way.

\medskip
\noindent\textbf{Underlying and market.}
\begin{itemize}
    \item $S$: price of the underlying asset (the spot), rescaled so that each trading day opens at $1$.
    \item $v$: instantaneous variance of the underlying; the volatility is $\sqrt{v}$.
    \item $\sigma$: volatility, with $\sigma = \sqrt{v}$; used in the Greek definitions of Section~1.
    \item $t$: calendar time. Written as a prefix, $\Delta t$ is a single hedging interval (one five-minute bar in the test).
    \item $\mu$: real-world (physical) drift of the underlying.
    \item $r$: risk-free interest rate ($r = 0.045$ in the test).
    \item $X,\,W$: the Brownian motions driving the price and the variance, correlated by $dX\,dW = \rho\,dt$.
\end{itemize}

\noindent\textbf{Heston parameters.}
\begin{itemize}
    \item $\kappa$: rate at which the variance reverts toward its long-run level.
    \item $\eta$: long-run (mean) level of the variance.
    \item $\epsilon$: volatility of volatility.
    \item $\rho$: correlation between the price and variance shocks.
    \item $\lambda$: market price of volatility risk; $\lambda = 0$ in the test, so the physical and risk-neutral variance drifts coincide.
    \item $\nu_{\mathrm{Q}}$: risk-neutral drift of the variance, $\nu_{\mathrm{Q}} = \kappa(\eta - v) - \lambda v$ (distinct from Vega $\nu$).
\end{itemize}

\noindent\textbf{Option and hedge portfolio.}
\begin{itemize}
    \item $V$: value of the option being hedged (a fresh at-the-money 0DTE call).
    \item $V_1$: value of the second option held to hedge variance (Vega) risk.
    \item $\Pi$: value of the delta--vega hedged portfolio, $\Pi = V - \Delta S - \Delta_1 V_1$.
    \item $\Delta$: units of the underlying held in the hedge, $\Delta = \partial V/\partial S - \Delta_1\,\partial V_1/\partial S$. Written as a prefix ($\Delta S$, $\Delta v$, $\Delta\Pi$) it instead denotes the change in a quantity over one hedging interval.
    \item $\Delta_1$: units of the hedging option $V_1$ held, chosen to cancel the portfolio's net Vega, $\Delta_1 = (\partial V/\partial v)\big/(\partial V_1/\partial v)$.
\end{itemize}

\begin{thebibliography}{99}

\bibitem{blackscholes1973}
Black, F. and Scholes, M. (1973).
\textit{The Pricing of Options and Corporate Liabilities.}
Journal of Political Economy, 81(3), 637--654.

\bibitem{merton1973}
Merton, R. C. (1973).
\textit{Theory of Rational Option Pricing.}
The Bell Journal of Economics and Management Science, 4(1), 141--183.

\bibitem{heston1993}
Heston, S. L. (1993).
\textit{A Closed-Form Solution for Options with Stochastic Volatility with Applications to Bond and Currency Options.}
The Review of Financial Studies, 6(2), 327--343.

\bibitem{gatheral2006}
Gatheral, J. (2006).
\textit{The Volatility Surface: A Practitioner's Guide.}
Wiley, Hoboken, NJ.

\bibitem{elkaroui1998}
El~Karoui, N., Jeanblanc-Picqu\'e, M., and Shreve, S. E. (1998).
\textit{Robustness of the Black and Scholes Formula.}
Mathematical Finance, 8(2), 93--126.

\bibitem{boyleemanuel1980}
Boyle, P. P. and Emanuel, D. (1980).
\textit{Discretely Adjusted Option Hedges.}
Journal of Financial Economics, 8(3), 259--282.

\bibitem{kloedenplaten1992}
Kloeden, P. E. and Platen, E. (1992).
\textit{Numerical Solution of Stochastic Differential Equations.}
Springer-Verlag, Berlin.

\bibitem{lord2010}
Lord, R., Koekkoek, R., and van Dijk, D. (2010).
\textit{A Comparison of Biased Simulation Schemes for Stochastic Volatility Models.}
Quantitative Finance, 10(2), 177--194.

\bibitem{kamal1999}
Kamal, M. (1999).
\textit{When You Cannot Hedge Continuously: The Corrections to Black--Scholes.}
Risk, 12(1), 82--85.

\bibitem{bertsimas2000}
Bertsimas, D., Kogan, L., and Lo, A. W. (2000).
\textit{When Is Time Continuous?}
Journal of Financial Economics, 55(2), 173--204.

\bibitem{bergomi2016}
Bergomi, L. (2016).
\textit{Stochastic Volatility Modeling.}
Chapman and Hall/CRC, Boca Raton, FL.

\bibitem{dimeraker2023}
Dim, C., Eraker, B., and Vilkov, G. (2023).
\textit{0DTEs: Trading, Gamma Risk and Volatility Propagation.}
Working paper, SSRN 4692190.

\bibitem{bozovic2025}
Bo\v{z}ovi\'c, M. (2025).
\textit{Intraday Jumps and 0DTE Options: Pricing and Hedging Implications.}
Working paper, SSRN 5223127.

\bibitem{cboe2023vix1d}
Cboe Global Markets (2023).
\textit{Cboe 1-Day Volatility Index (VIX1D) Methodology.}
Cboe Global Markets.

\bibitem{Maruyama1955}
Maruyama, G. (1955).
\textit{Continuous Markov Processes and Stochastic Equations.}
Rendiconti del Circolo Matematico di Palermo, 4(1), 48--90.

\end{thebibliography}

\end{document}
